%!TEX root = ../../main.tex
\section{관측 갱신 층화}
\label{sec:val-frame}

Match-V5의 상태는 1분 프레임으로만 갱신되므로, 예측 시점과 결과 시점 사이에 새 프레임이 들어왔는지와 예측 시점의 프레임이 얼마나 오래된 것인지는 행마다 다르다. 이 두 조건으로 15.16 TEST를 2$\times$2로 층화하고, 동결 $q$와 $\PTflex$의 점수를 셀 안에서 경기 가중으로 다시 계산하였다(표 \ref{tab:frame-strata}; 보완 실험, 15.16 노출 이후의 진단). 셀은 같은 프레임 / 새 프레임 $\times$ 예측 시점 프레임 나이 30초 미만 / 이상이다. 예측 시점의 프레임 나이 중앙값은 T \fsTPreAgeMed{}초, S \fsSPreAgeMed{}초이고, 결과 시점까지 새 프레임이 없는 행은 T의 \fsTSameFrame{}\%, S의 \fsSSameFrame{}\%이다.

%% GENERATED by thesis/tools/gen_supp.py -- do not edit
\begin{table}[!t]
  \centering
  \caption{관측 갱신 층화별 결과 (15.16 TEST)}
  \label{tab:frame-strata}
  \footnotesize

  \begin{tabular}{@{}llrrrrrrr@{}}
    \toprule
    코호트 & 셀 & $n$ & 경기 & $P(\mathrm{SVI}{=}1)$ & $\mathbb{E}|\Delta\hat V|$ & $\Delta$Brier & $\bar p_{\mathrm{pre}}$ & $\bar t$ (분) \\
    \midrule
    T & 같은 프레임, 나이 $<$30 s & 2{,}204 & 2{,}164 & 0.530 & 0.094 & -0.00090 & 0.497 & 22.9 \\
    T & 같은 프레임, 나이 $\ge$30 s & 20 & 20 & 0.450 & 0.057 & +0.00427 & 0.498 & 16.6 \\
    T & 새 프레임, 나이 $<$30 s & 10{,}003 & 8{,}979 & 0.486 & 0.119 & -0.00347 & 0.504 & 21.2 \\
    T & 새 프레임, 나이 $\ge$30 s & 20{,}754 & 16{,}923 & 0.497 & 0.116 & -0.00398 & 0.506 & 21.1 \\
    T & 전체 & 32{,}981 & 24{,}020 & 0.496 & 0.115 & -0.00373 & 0.505 & 21.3 \\
    S & 같은 프레임, 나이 $<$30 s & 12{,}912 & 11{,}712 & 0.569 & 0.078 & -0.00141 & 0.498 & 10.7 \\
    S & 같은 프레임, 나이 $\ge$30 s & 1{,}114 & 1{,}101 & 0.590 & 0.077 & +0.00007 & 0.504 & 9.2 \\
    S & 새 프레임, 나이 $<$30 s & 31{,}314 & 24{,}695 & 0.493 & 0.091 & -0.00396 & 0.501 & 11.1 \\
    S & 새 프레임, 나이 $\ge$30 s & 55{,}865 & 36{,}561 & 0.503 & 0.091 & -0.00310 & 0.506 & 12.5 \\
    S & 전체 & 101{,}205 & 49{,}730 & 0.509 & 0.089 & -0.00335 & 0.504 & 11.8 \\
    \bottomrule
  \end{tabular}
  \tabnote{셀 = 예측 시점과 결과 시점 사이에 새 분 프레임이 들어왔는지(같은/새 프레임) $\times$ 예측 시점 프레임의 나이(30 s 기준). $\Delta$Brier = $q - \mathrm{PT}_{\mathrm{flex}}$, 동결 점수, 셀 안에서 경기 가중 재계산, 구간 없음. 조성의 진단이며 인과 효과가 아니다.}
\end{table}


새 프레임이 들어온 두 셀에서 $\dBrier$는 T에서 \fsTnyDb{}와 \fsTnoDb{}, S에서 \fsSnyDb{}와 \fsSnoDb{}로 전체 값과 같은 부호이다. 같은 프레임 셀에서는 값이 0에 가깝거나 부호가 바뀌며(T 같은 프레임·30초 이상 셀은 \fsTsoN{}행에 불과하고 \fsTsoDb{}; S 같은 프레임·30초 미만 셀은 \fsSsyDb{}), 같은 프레임 셀은 $\SVI$ 양성률과 $\E|\dV|$의 분포도 다르다(T 같은 프레임·30초 미만 셀의 양성률 \fsTsyPsvi{}, $\E|\dV|$ \fsTsyEdv{}). 이 차이는 동결 점수 아래에서 셀의 조성이 다르다는 사실의 기술이며, 프레임 갱신이 성능을 만들어 냈다는 인과 문장의 근거가 아니다. 결과 시점이나 프레임 갱신 여부는 예측 시점 이후의 정보이므로 $q$의 입력에 들어가지 않는다.
